\chapter{\eh{rocket} Introducción a la Parte 2 --- Nivel Senior}

%% ── 0.1 ¿Qué cubre esta parte? ──────────────────────────────────────────────
\section{\eh{books} ¿Qué cubre esta parte?}

\begin{concepto}
La \textbf{Parte 1} de estos apuntes cubrió los \textbf{fundamentos de redes}:
el modelo OSI, TCP/IP, protocolos clave (TCP, UDP, DNS, HTTP, TLS, SSH) y
cómo fluyen los datos de una máquina a otra.

La \textbf{Parte 2} sube un nivel. Aquí vemos cómo funcionan las redes
\textbf{por dentro}, en producción real, a escala empresarial y de nube:

\begin{itemize}
    \item \textbf{Datacenters modernos:} VxLAN, EVPN, Spine-Leaf, SDN
    \item \textbf{Redes de operador (carrier):} MPLS, QoS, ingeniería de tráfico
    \item \textbf{Kubernetes:} CNI, Service Mesh, Ingress, eBPF
    \item \textbf{Cloud networking:} AWS VPC, Transit Gateway, Direct Connect
    \item \textbf{Seguridad de switch empresarial:} 802.1X, MACsec, VLAN hardening
    \item \textbf{Wi-Fi empresarial:} WPA3-Enterprise, 802.11ax (Wi-Fi 6)
    \item \textbf{Multicast:} PIM, IGMP, streaming a gran escala
\end{itemize}

\emoji{light-bulb} \textbf{Analogía:}
La Parte 1 es aprender a \textbf{manejar un auto}: sabes presionar el acelerador,
girar el volante, respetar semáforos.
La Parte 2 es aprender a \textbf{construir el motor}: sabes cómo funciona
el sistema de combustión, la transmisión, el turbocompresor --- y por qué
falla a las 3am cuando hay 50{,}000 usuarios conectados.
\end{concepto}

%% ── 0.2 El Stack del Datacenter Moderno ─────────────────────────────────────
\section{\eh{office-building} El Stack del Datacenter Moderno}

Antes de entrar a cada tecnología, necesitas ver el panorama completo.
Este es el flujo real de un paquete desde Internet hasta el contenedor
que lo procesa en un datacenter moderno (AWS, Google Cloud, un banco grande):

\begin{center}
\begin{lstlisting}
Internet
   |
   v
Edge Routers  (BGP — hablan con todo Internet)
   |
   v
Firewall / DDoS Mitigation  (filtra ataques, limita flujos)
   |
   v
Load Balancers  (L4/L7 — distribuyen tráfico entre servidores)
   |
   v
Spine Switches  (ECMP — múltiples rutas iguales, sin STP)
   |
   v
Leaf Switches  (VxLAN VTEP — virtualizan la red sobre IP)
   |
   v
Servers / Hypervisors  (Linux + Open vSwitch o eBPF)
   |
   v
Pods / Containers  (Kubernetes — la app vive aquí)
\end{lstlisting}
\end{center}

\begin{concepto}
Cada flecha del diagrama es un capítulo de esta parte:

\begin{itemize}
    \item \textbf{BGP} se estudia en Cloud Networking (cap.\ 07).
    \item \textbf{Firewall / DDoS} es parte de Seguridad empresarial (cap.\ 02).
    \item \textbf{Load Balancers + VxLAN + Spine-Leaf} son el cap.\ 04 (Datacenter).
    \item \textbf{MPLS y QoS} mueven paquetes entre datacenters (cap.\ 05).
    \item \textbf{Pods y Kubernetes} son el cap.\ 06.
\end{itemize}

\emoji{light-bulb} \textbf{Analogía:}
El datacenter es como un aeropuerto enorme. Internet es el mundo exterior.
Los Edge Routers son la \textit{aduana}. Los Spine Switches son las
\textit{pistas de aterrizaje} (no hay semáforos --- todo va en paralelo).
Los Leaf Switches son las \textit{puertas de abordaje}. Los Pods son
\textit{los asientos de cada vuelo}: el destino final del pasajero (paquete).
\end{concepto}

%% ── 0.3 Mapa de Tecnologías ─────────────────────────────────────────────────
\section{\eh{world-map} Mapa de Tecnologías}

La siguiente tabla muestra qué cubre cada capítulo, qué tecnologías
involucra y en qué entorno vive en producción real:

\begin{center}
\renewcommand{\arraystretch}{1.4}
\begin{tabular}{>{\bfseries}c p{4.5cm} p{5.5cm}}
\toprule
\textbf{Cap.} & \textbf{Tecnología} & \textbf{Dónde vive} \\
\midrule
01 & Wi-Fi Empresarial: 802.11ax, WPA3-Enterprise, RADIUS, roaming 802.11r
   & Oficinas corporativas, campus universitarios, hoteles, aeropuertos \\
02 & Switching Avanzado: VLANs, STP/RSTP, 802.1X, MACsec, port security
   & Cualquier red con switches administrables (empresa, datacenter, banco) \\
03 & SDN y OpenFlow: controladores, APIs northbound/southbound, OpenDaylight
   & Datacenters de hiperescaladores (Google, Facebook), redes universitarias \\
04 & Datacenter: VxLAN, EVPN, Spine-Leaf, ECMP, BGP-EVPN
   & AWS, Azure, GCP, Equinix, datacenters privados grandes \\
05 & MPLS y QoS: LSP, LDP, RSVP-TE, ingeniería de tráfico, DiffServ
   & Redes de operador (Telmex, AT\&T, Lumen), troncales entre datacenters \\
06 & Kubernetes Networking: CNI, Calico, Cilium/eBPF, Service Mesh, Ingress
   & Toda empresa que corre microservicios en producción \\
07 & Cloud Networking: VPC, Transit Gateway, Direct Connect, BGP con AWS
   & AWS, Azure, GCP --- el 90\% de las startups y empresas medianas \\
08 & Multicast: PIM-SM, PIM-SSM, IGMP, MDT, streaming y finanzas
   & Broadcasters (televisión IP), bolsas de valores, CDNs \\
\bottomrule
\end{tabular}
\end{center}

%% ── 0.4 Parte 1 vs Parte 2 ──────────────────────────────────────────────────
\section{\eh{balance-scale} Parte 1 vs Parte 2 --- ¿Qué cambia?}

Esta tabla compara el nivel de cada tema entre las dos partes.
No es que Parte 1 sea ``incorrecto'' --- es que Parte 2 es lo que
se usa en producción a escala:

\begin{center}
\renewcommand{\arraystretch}{1.4}
\begin{tabular}{p{3.8cm} p{4cm} p{4cm}}
\toprule
\textbf{Dimensión} & \textbf{Parte 1 (Fundamentos)} & \textbf{Parte 2 (Producción)} \\
\midrule
Protocolos de transporte
   & TCP / UDP básico, puertos, handshake
   & MPLS / VxLAN, túneles, encapsulación multi-capa \\
Seguridad
   & TLS (HTTPS), SSH, firewalls básicos
   & MACsec (capa 2), 802.1X, WPA3-Enterprise, Zero Trust \\
Wi-Fi
   & 802.11 básico, WPA2 personal, canales
   & Wi-Fi 6 (802.11ax), roaming empresarial, RADIUS/EAP \\
Enrutamiento
   & RIP, OSPF básico, tablas de rutas
   & BGP completo, MPLS-TE, ECMP en Spine-Leaf \\
Switches
   & VLANs básicas, STP para evitar loops
   & EVPN, VxLAN VTEP, MLAG, port security avanzado \\
Virtualización de red
   & No aplica (redes físicas)
   & SDN, OpenFlow, OVS, eBPF, Cilium \\
Contenedores
   & No aplica
   & CNI plugins, Kubernetes networking, Service Mesh \\
Cloud
   & No aplica
   & VPC, Transit Gateway, Direct Connect, Cloud Router \\
Escala
   & Red de laboratorio o SOHO (10--100 hosts)
   & Miles de servidores, millones de flujos simultáneos \\
\bottomrule
\end{tabular}
\end{center}

%% ── 0.5 Por qué importa ─────────────────────────────────────────────────────
\section{\eh{graduation-cap} ¿Por qué importa todo esto?}

\begin{moderno}
\textbf{El 90\% de la infraestructura empresarial moderna corre sobre:}

\begin{itemize}
    \item \textbf{Kubernetes} para orquestar contenedores
    \item \textbf{Cloud networking} (AWS VPC, Azure VNet) para conectar todo
    \item \textbf{MPLS} en las redes de operador que unen datacenters
    \item \textbf{VxLAN + EVPN} dentro del datacenter para virtualizar la red
    \item \textbf{802.1X + MACsec} para seguridad de acceso en campus empresariales
\end{itemize}

\vspace{0.5em}
\textbf{Las entrevistas de SRE / Network Engineer en empresas como Google,
Amazon, Banco de México, BBVA, Telmex y cualquier empresa con más de
500 empleados preguntan exactamente estos temas.}

\vspace{0.5em}
\emoji{brain} Si sabes explicar por qué VxLAN usa UDP port 4789,
cómo Kubernetes asigna IPs a los pods, y qué hace un LSR en MPLS ---
ya estás en el top 10\% de candidatos.

\vspace{0.5em}
\emoji{light-bulb} \textbf{Analogía final:}
Un médico general sabe que el corazón bombea sangre (Parte 1).
Un cirujano cardiovascular sabe hacer un bypass, reparar una válvula
y gestionar la UCI post-operatoria (Parte 2).
Ambos son necesarios, pero el hospital solo le confía el quirófano al segundo.
\end{moderno}
